\section{Methodology: Hessian-Regularized Coefficient Optimization}
\label{sec:methodology}

We present \textbf{Hessian-Regularized Coefficient Optimization (HessMerge)}, a mathematically rigorous framework designed to stabilize test-time adaptive model merging under downstream deployment quantization. HessMerge operates without requiring the original multi-task training datasets, optimizing layer-wise continuous coefficients purely on a lightweight calibration stream of unlabeled test-set queries.

\subsection{Differentiable Layer-Wise Parameterization}
Let $\mathcal{M}_{\text{pre}}$ denote a pre-trained base model (e.g., a Vision Transformer backbone), and let $\{\mathcal{M}_k\}_{k=1}^K$ represent a set of $K$ independent task-specific expert models fine-tuned from the same shared pre-trained initialization. We denote their respective parameter weight tensors at layer group $l \in \{1, \dots, L\}$ as $W_{\text{pre}}^l$ and $W_k^l$. 

The task vector for expert $k$ at layer $l$ is defined as the parameter difference:
\begin{equation}
    \tau_k^l = W_k^l - W_{\text{pre}}^l.
\end{equation}

Rather than applying a single, static coefficient across the entire network, HessMerge parameterizes the merged weight tensor $W_{\text{merged}}^l(\boldsymbol{\lambda}^l)$ at each layer group $l$ dynamically using layer-wise blending coefficients $\boldsymbol{\lambda}^l = [\lambda_1^l, \dots, \lambda_K^l]^T \in [0, 1]^K$:
\begin{equation}
    W_{\text{merged}}^l(\boldsymbol{\lambda}^l) = W_{\text{pre}}^l + \sum_{k=1}^K \lambda_k^l \tau_k^l.
\end{equation}

To enforce the hard box constraint $\lambda_k^l \in [0, 1]$ during gradient-based optimization and maintain a smooth, twice-differentiable mapping, we parameterize the raw coefficients using real-valued variables $\theta_k^l \in \mathbb{R}$ mapped through the logistic sigmoid function:
\begin{equation}
    \lambda_k^l = \sigma(\theta_k^l) = \frac{1}{1 + e^{-\theta_k^l}}.
\end{equation}
The full optimization matrix is denoted by $\Theta \in \mathbb{R}^{L \times K}$.

\subsection{Quantization Noise and Loss Landscape Curvature}
When the continuous merged parameters are deployed onto edge hardware, they are subjected to a post-training quantization operator $\mathcal{Q}(\cdot)$ (e.g., INT8 symmetric per-channel or INT4 symmetric per-channel). The quantized parameter weights can be represented as:
\begin{equation}
    W_{\text{quant}}^l = \mathcal{Q}\left(W_{\text{merged}}^l(\boldsymbol{\lambda}^l)\right) = W_{\text{merged}}^l(\boldsymbol{\lambda}^l) + \delta^l,
\end{equation}
where $\delta^l$ represents the multi-dimensional quantization noise tensor at layer $l$.

To analyze the impact of quantization noise on task performance, we perform a second-order Taylor expansion of the multi-task loss function $\mathcal{L}$ around the continuous optimized parameter state $W_{\text{merged}}$:
\begin{equation}
    \mathcal{L}(W_{\text{quant}}) \approx \mathcal{L}(W_{\text{merged}}) + \nabla_W \mathcal{L}(W_{\text{merged}})^T \delta + \frac{1}{2} \delta^T \nabla^2_W \mathcal{L}(W_{\text{merged}}) \delta.
\end{equation}

Assuming the optimization of the blending coefficients has successfully converged to a local minimum under the continuous objective, the first-order gradient term vanishes ($\nabla_W \mathcal{L}(W_{\text{merged}}) \approx 0$). The performance drop due to quantization is therefore dominated by the quadratic curvature term:
\begin{equation}
    \Delta \mathcal{L}_{\text{quant}} = \mathcal{L}(W_{\text{quant}}) - \mathcal{L}(W_{\text{merged}}) \approx \frac{1}{2} \delta^T \mathcal{H}_W \delta,
\end{equation}
where $\mathcal{H}_W = \nabla^2_W \mathcal{L}(W_{\text{merged}})$ is the parameter-space Hessian matrix. 

Under standard PTQ schemes, the maximum rounding error is bounded by half the quantization step size ($s/2$). Consequently, the norm of the quantization noise $\|\delta\|_2^2$ is bounded. The critical factor determining the loss gap is the curvature of the loss landscape. By Cauchy-Schwarz, the loss gap is upper-bounded by:
\begin{equation}
    \Delta \mathcal{L}_{\text{quant}} \le \frac{1}{2} \lambda_{\max}(\mathcal{H}_W) \|\delta\|_2^2,
\end{equation}
where $\lambda_{\max}(\mathcal{H}_W)$ is the maximum eigenvalue of the Hessian, representing the sharpness of the local loss minimum.

\subsection{The HessMerge Objective: Hessian Trace Regularization}
Directly computing or regularizing the parameter-space Hessian $\mathcal{H}_W$ is computationally intractable due to the massive dimensionality of modern neural networks. However, because the blending coefficients $\Theta$ are extremely low-dimensional ($L \times K \approx 56$), we can project the loss curvature directly into the coefficient space $\Theta$.

HessMerge optimizes the coefficient matrix $\Theta$ to minimize the joint objective:
\begin{equation}
    \mathcal{L}_{\text{HessMerge}}(\Theta) = \mathcal{L}_{\text{entropy}}(\Theta) + \gamma \cdot \text{Tr}\left(\mathcal{H}_{\Theta}\right),
\end{equation}
where $\mathcal{L}_{\text{entropy}}$ is the unsupervised multi-task entropy loss, $\gamma > 0$ is the regularization strength, and $\text{Tr}\left(\mathcal{H}_{\Theta}\right)$ is the trace of the Hessian with respect to the raw coefficients $\Theta$:
\begin{equation}
    \text{Tr}\left(\mathcal{H}_{\Theta}\right) = \sum_{l=1}^L \sum_{k=1}^K \frac{\partial^2 \mathcal{L}_{\text{entropy}}(\Theta)}{\partial (\theta_k^l)^2}.
\end{equation}

Minimizing the Hessian trace explicitly flattens the local loss landscape with respect to the merging coefficients. This ensures that any perturbation in the parameter space (induced by quantization) results in a minimal performance drop, guaranteeing robustness across diverse deployment platforms.
